\uuid{dtGv}
\exo7id{2736}
\titre{exo7 2736}
\auteur{tumpach}
\organisation{exo7}
\datecreate{2009-10-25}
\isIndication{false}
\isCorrection{false}
\chapitre{Déterminant, système linéaire}
\sousChapitre{Système linéaire, rang}
\cpge{mpsi-mp2i.calcul-matriciel-systemes.c, pcsi.calcul-matriciel-systemes-lineaires.c, ptsi.calcul-matriciel-systemes-lineaires.c, mpsi-mp2i.calcul-matriciel-systemes.a, pcsi.calcul-matriciel-systemes-lineaires.a, ptsi.calcul-matriciel-systemes-lineaires.a, mpsi-mp2i.applications-lineaires.a, pcsi.applications-lineaires.a, ptsi.applications-lineaires.a, bcpst1.systemes-lineaires.a, bcpst1.matrices, bcpst1.applications-lineaires-matrices}
\module{Algèbre}
\niveau{L2}
\difficulte{}

\contenu{
\texte{
Les vecteurs complexes $(z, w)$ et $(z', w')$ sont li\'es par la formule $(z', w') = (z + iw, (1+i)z + (1-2i)w)$. Un \'etudiant qui n'aime pas les nombres complexes pose $z = x + iy$, $w = u + iv$, $z' = x' + i y'$ et $w' = u'+iv'$.
}
\begin{enumerate}
    \item \question{Exprimer $(x', y', u', v')$ en fonction de $(x, y, u, v)$.}
    \item \question{R\'esoudre le syst\`eme $(x', y', u', v') = (1, 2, 3, 4)$.}
\end{enumerate}
}
